% sections/future_work.tex --- Future Work
\section{Future Work}
\label{sec:future-work}

We close with a list of extensions that the present release does
not cover.  Each is supported by a configuration switch in
\system{}, an export already produced by the runtime, or a related
artefact in the repository, and so represents a relatively
low-friction path from v$1.0$ of \bench{} to a richer benchmark.

\subsection{Multi-period and rolling-window benchmarks}
\label{sec:future-work:rolling}

The v$1.0$ temporal split holds out the last three fiscal months
of a single year.  A more demanding regime is the rolling-window
protocol that mirrors continuous-auditing practice: the model is
re-trained at the close of each month on all data up to that
point and evaluated on the following month, repeated across the
year.  This regime exposes whether a method's accuracy degrades
under \emph{distribution shift within} the period---a more
realistic deployment scenario than a single hold-out.  The
required configuration changes are local to the splits files; no
new generator features are needed.

\subsection{Cross-jurisdictional comparability}
\label{sec:future-work:jurisdiction}

\system{}'s ISO~21378 alignment (Section~\ref{sec:datasynth:iso})
populates account-class and account-sub-class fields that are
common across jurisdictions but does so on a single chart of
accounts at a time.  A natural next step is to construct splits
that hold out entire jurisdictions---training on US-GAAP and
IFRS, evaluating on French- and German-GAAP charts---and to
report the resulting drop in performance as a measure of
cross-jurisdictional transfer.  This is exactly the kind of
question the standardised account-class fields were designed to
support.

\subsection{Counterfactual evaluation}
\label{sec:future-work:counterfactual}

The \system{} counterfactual engine produces, for any base
generation, a paired \emph{counterfactual} dataset in which a
specified perturbation (an anomaly removed, a different control in
place, a different subsidiary structure) propagates through the
causal DAG.  Pairing \bench{} edges with their counterfactuals
yields a strictly stronger evaluation: a method should not only
flag the anomalous edge in the base run, but also \emph{not} flag
the corresponding edge in the counterfactual run where the
anomaly is absent.  We anticipate releasing counterfactual pairs
in v$1.1$.

\subsection{Real-data calibration loop}
\label{sec:future-work:calibration}

The empirical foundations of the generator
(Section~\ref{sec:empirical}) were established on a $2024$
corpus.  The corpus is not static---new client engagements
continually produce additional ledgers---and so the calibration
of the generator should be revisited periodically.  We propose
to publish, with each minor version of \bench{}, a refresh of the
empirical Tables~\ref{tab:emp:size}--\ref{tab:emp:bin1} and a
delta against the previous calibration; benchmark numbers are
versioned to specific calibrations so that changes are
traceable.

\subsection{Beyond anomaly detection}
\label{sec:future-work:beyond}

The benchmark schema is a general-purpose accounting graph and
supports tasks adjacent to anomaly detection: control testing
(does the method correctly identify the controls implicated by a
given posting?), account-relationship discovery (does the method
recover the canonical chart-of-accounts hierarchy from the
edge list?), and cutoff testing (does the method correctly
classify postings near period boundaries as in-period or
out-of-period?).  These tasks would be released as separate
``challenges'' on the same underlying graph---a model trained
for anomaly detection could be evaluated against them
zero-shot, and a model trained jointly across challenges might
benefit from the additional supervision.
